\section{Metodología}

\subsection{Dinámica histórica bajo la medida física}

Se supone que el activo sigue un movimiento browniano geométrico bajo la medida física:
\begin{equation}
 dS_t=\mu S_t\,dt+\sigma S_t\,dW_t^{\Pmeasure}.
 \label{eq:p_gbm}
\end{equation}
Para observaciones igualmente espaciadas con intervalo $\Delta t$, los log-retornos
\begin{equation}
 R_i=\log\left(\frac{S_{t_i}}{S_{t_{i-1}}}\right)
\end{equation}
satisfacen
\begin{equation}
 R_i\mid\mu,\sigma
 \sim
 \mathcal N\!\left[
 \left(\mu-\frac12\sigma^2\right)\Delta t,
 \sigma^2\Delta t
 \right].
 \label{eq:return_likelihood}
\end{equation}
La función de verosimilitud se obtiene como el producto de estas densidades condicionales.

\subsection{Especificación bayesiana y Metropolis--Hastings}

La especificación base conserva los priors del proyecto original:
\begin{align}
 \mu &\sim \mathcal N(0,1),\\
 \sigma &\sim \operatorname{InvGamma}(2,0.1), \qquad \sigma>0,
 \label{eq:priors}
\end{align}
donde la segunda densidad es proporcional a $\sigma^{-3}\exp(-0.1/\sigma)$. Esta parametrización se mantiene como baseline y no se interpreta como una elección universal; la sensibilidad al prior se discute como una extensión pendiente.

Para respetar automáticamente la restricción $\sigma>0$, el algoritmo trabaja con
\begin{equation}
 \eta=\log\sigma.
\end{equation}
La densidad objetivo en $(\mu,\eta)$ incorpora el Jacobiano de la transformación:
\begin{equation}
 \log\pi(\mu,\eta\mid\mathcal D)
 =
 \log L(\mu,e^\eta;\mathcal D)
 +\log p(\mu)
 +\log p(e^\eta)
 +\eta + C.
 \label{eq:posterior_eta}
\end{equation}
Se utiliza una propuesta Random-Walk Metropolis simétrica. Si $\theta=(\mu,\eta)$, se propone
\begin{equation}
 \theta'=\theta+\varepsilon,
 \qquad
 \varepsilon\sim\mathcal N(0,\Sigma_q),
\end{equation}
y la probabilidad de aceptación es
\begin{equation}
 \alpha(\theta,\theta')
 =\min\left\{1,
 \frac{\pi(\theta'\mid\mathcal D)}{\pi(\theta\mid\mathcal D)}
 \right\}.
\end{equation}
La implementación sigue el principio general de \citet{metropolis1953} y \citet{hastings1970}.

\subsection{Valuación bajo la medida neutral al riesgo}

En Black--Scholes--Merton con rendimiento continuo por dividendos $q$, la dinámica de valuación es
\begin{equation}
 dS_t=(r-q)S_t\,dt+\sigma S_t\,dW_t^{\Qmeasure}.
 \label{eq:q_gbm}
\end{equation}
El drift físico $\mu$ no aparece en \eqref{eq:q_gbm}. Por ello, aunque $p(\mu,\sigma\mid\mathcal D)$ sea conjunta, la incertidumbre de $\mu$ no se propaga al precio de no arbitraje dentro de esta especificación.

Sea una call asiática aritmética discretamente monitoreada en $0<t_1<\cdots<t_m=T$ con
\begin{equation}
 A_T=\frac1m\sum_{j=1}^{m}S_{t_j}.
\end{equation}
Condicionado en $\sigma$, el precio es
\begin{equation}
 C^{\Qmeasure}(\sigma)
 =e^{-rT}\E^{\Qmeasure}_{\sigma}
 \left[(A_T-K)^+\right].
 \label{eq:cq}
\end{equation}
La posterior de volatilidad induce entonces una distribución posterior de precios por push-forward:
\begin{equation}
 p(C\mid\mathcal D)
 =\int
 \delta_{C^{\Qmeasure}(\sigma)}(C)
 p(\sigma\mid\mathcal D)\,d\sigma.
 \label{eq:pushforward}
\end{equation}
Esta distribución representa incertidumbre paramétrica sobre un precio teórico; no es una distribución de pérdidas de cartera ni una medida de VaR o Expected Shortfall.

\subsection{Monte Carlo y reducción de varianza}

El precio aritmético se aproxima mediante simulación exacta de la GBM en las fechas de monitoreo. Se utilizan pares antitéticos y una opción asiática geométrica como variable de control. Sea
\begin{equation}
 G_T=\left(\prod_{j=1}^{m}S_{t_j}\right)^{1/m}.
\end{equation}
Si $Y=e^{-rT}(A_T-K)^+$ y $X=e^{-rT}(G_T-K)^+$, con $\E[X]$ disponible analíticamente, el estimador de control es
\begin{equation}
 Y_{cv}=Y-\beta\left(X-\E[X]\right),
 \qquad
 \beta=\frac{\operatorname{Cov}(Y,X)}{\Var(X)}.
 \label{eq:cv}
\end{equation}
La estrategia sigue a \citet{kemnavorst1990}. En las mallas de pricing se emplean además números aleatorios comunes a través de los puntos de volatilidad de un mismo contrato, con el objetivo de suavizar comparaciones e interpolación.

\subsection{Reglas de valuación comparadas}

Se consideran cuatro reglas:
\begin{align}
 \widehat C_{FB}
 &=\E\left[C^{\Qmeasure}(\sigma)\mid\mathcal D\right],
 \label{eq:fb}\\
 \widehat C_{PM}
 &=C^{\Qmeasure}\!\left(\E[\sigma\mid\mathcal D]\right),
 \label{eq:pm}\\
 \widehat C_{MAP}
 &=C^{\Qmeasure}(\widehat\sigma_{MAP}),
 \label{eq:map}\\
 \widehat C_{MLE}
 &=C^{\Qmeasure}(\widehat\sigma_{MLE}).
 \label{eq:mle}
\end{align}

En general, $\widehat C_{FB}$ y $\widehat C_{PM}$ no tienen por qué coincidir. Una expansión de Taylor de segundo orden alrededor de $\bar\sigma=\E[\sigma\mid\mathcal D]$ produce
\begin{equation}
 \widehat C_{FB}
 \approx
 C^{\Qmeasure}(\bar\sigma)
 +\frac{1}{2}C^{\Qmeasure\prime\prime}(\bar\sigma)
 \Var(\sigma\mid\mathcal D).
 \label{eq:taylor}
\end{equation}
La ecuación \eqref{eq:taylor} sugiere dos predicciones. La brecha entre full Bayes y plug-in debe disminuir cuando la posterior de $\sigma$ se concentra, y debe aumentar en regiones donde el precio sea más curvo respecto de la volatilidad. El diseño experimental se construye para evaluar ambas predicciones.
